\documentclass[12pt,a4paper]{article}
\usepackage[utf8]{inputenc}
\usepackage[T1]{fontenc}
\usepackage[spanish]{babel}
\usepackage{geometry}
\geometry{a4paper, top=2.5cm, bottom=2.5cm, left=2.5cm, right=2.5cm}
\usepackage{listings}
\usepackage{xcolor}
\usepackage{tcolorbox}
\usepackage{hyperref}
\usepackage{graphicx}
\usepackage{tikz}
\usetikzlibrary{positioning, arrows.meta}

\definecolor{codegreen}{rgb}{0,0.6,0}
\definecolor{codegray}{rgb}{0.5,0.5,0.5}
\definecolor{codepurple}{rgb}{0.58,0,0.82}
\definecolor{backcolour}{rgb}{0.96,0.96,0.96}

\lstdefinestyle{javastyle}{
    language=Java,
    backgroundcolor=\color{backcolour},
    commentstyle=\color{codegreen}\itshape,
    keywordstyle=\color{blue}\bfseries,
    numberstyle=\tiny\color{codegray},
    stringstyle=\color{codepurple},
    basicstyle=\ttfamily\footnotesize,
    breakatwhitespace=false,
    breaklines=true,
    captionpos=t,
    keepspaces=true,
    numbers=left,
    numbersep=5pt,
    showspaces=false,
    showstringspaces=false,
    showtabs=false,
    tabsize=4,
    frame=single,
    rulecolor=\color{lightgray}
}
\lstset{
    style=javastyle,
    literate={á}{{\'a}}1 {é}{{\'e}}1 {í}{{\'i}}1 {ó}{{\'o}}1 {ú}{{\'u}}1
             {Á}{{\'A}}1 {É}{{\'E}}1 {Í}{{\'I}}1 {Ó}{{\'O}}1 {Ú}{{\'U}}1
             {ñ}{{\~n}}1 {Ñ}{{\~N}}1
}

\title{\textbf{Guía de Solución: Programación Orientada a Objetos} \\ \large Actividad 5: Interfaz Gráfica con Archivos (CRUD)}
\author{Cristian}
\date{}

\begin{document}
\maketitle

\section{Repositorio de GitHub}
El código fuente de este proyecto se encuentra alojado en GitHub. 
\newline
\textbf{Enlace del repositorio:} \url{https://github.com/cioranemil/poo-trabajo-3-}

\section{Enunciado: Actividad 5}
Desarrollo de una Aplicación con Interfaz Gráfica (GUI) conectada a un archivo de datos local. El sistema debe contener un formulario provisto de botones funcionales para ejecutar de manera exitosa el patrón de persistencia \textbf{CRUD} (Create, Read, Update y Delete) sobre un archivo de texto secuencial usando Java.

\section{Casos de Uso}
\begin{table}[h]
    \centering
    \begin{tabular}{|p{3cm}|p{12cm}|}
    \hline
    \textbf{Caso de Uso} & \textbf{Administrar Biblioteca Personal (CRUD)} \\ \hline
    \textbf{Actor} & Usuario / Lector \\ \hline
    \textbf{Descripción} & El usuario utiliza la interfaz para llevar el inventario de libros leídos, modificarlos o borrarlos, guardándolos en el disco duro. \\ \hline
    \textbf{Flujo Principal} & 
    1. El usuario inicia la aplicación y visualiza su lista actual de libros en una tabla (\textbf{Read}).\newline
    2. Ingresa un Título, Autor y un Estado en el formulario.\newline
    3. Al hacer clic en ``Añadir'', el libro se escribe en \texttt{biblioteca.txt} (\textbf{Create}).\newline
    4. El usuario selecciona un libro de la tabla y modifica su estado (ej. de ``Pendiente'' a ``Terminado''), haciendo clic en ``Actualizar'' (\textbf{Update}).\newline
    5. El usuario selecciona un libro y oprime ``Eliminar'', borrándolo de la tabla y del archivo (\textbf{Delete}). \\ \hline
    \textbf{Flujo Alternativo} & 
    Si se intenta añadir un libro sin título o autor, o si el archivo está en uso por otro programa, el sistema levanta un \texttt{JOptionPane} avisando al usuario y denegando la transacción para proteger los datos. \\ \hline
    \end{tabular}
    \caption{Descripción del Caso de Uso Principal (Actividad 5)}
\end{table}

\section{Diagrama de Clases y Explicación}
\begin{figure}[h]
    \centering
    \begin{tikzpicture}[
        node distance=3cm,
        umlbox/.style={draw, inner sep=0pt, align=left, font=\footnotesize}
    ]
    
    \node[umlbox] (gui) {
        \begin{tabular}{l}
        \multicolumn{1}{c}{\textbf{VentanaBiblioteca}} \\
        \hline
        - txtTitulo, txtAutor: JTextField \\
        - cbEstado: JComboBox \\
        - tablaLibros: JTable \\
        \hline
        <<constructor>> VentanaBiblioteca() \\
        - cargarTabla(): void \\
        - crear(): void \\
        - actualizar(): void \\
        - borrar(): void
        \end{tabular}
    };
    
    \node[umlbox, right=2cm of gui] (backend) {
        \begin{tabular}{l}
        \multicolumn{1}{c}{\textbf{GestorArchivoBiblioteca}} \\
        \hline
        - FILE\_NAME: String \\
        \hline
        + leerLibros(): ArrayList \\
        + crearLibro(): void \\
        + actualizarLibro(): void \\
        + borrarLibro(): void
        \end{tabular}
    };
    
    \draw[-stealth, dashed] (gui) -- node[above] {Consume} (backend);
    
    \end{tikzpicture}
    \caption{Diagrama de clases del Módulo CRUD}
\end{figure}

\textbf{Explicación del diagrama de clases:}
La solución implementa una clara separación de responsabilidades (Modelo-Vista). La clase \texttt{VentanaBiblioteca} es estrictamente una capa de presentación gráfica, maneja los clics de botones y la actualización de la tabla visual (\texttt{JTable}). Por su parte, \texttt{GestorArchivoBiblioteca} contiene la lógica de persistencia; encapsula por completo a la clase \texttt{RandomAccessFile} del API de Java, asegurando que las operaciones de lectura, escritura y renombrado de archivos temporales sean opacas para la interfaz de usuario.

\section{Diagrama de Objetos}
\begin{verbatim}
[ uiBiblioteca: VentanaBiblioteca ] -.usa.-> [ backend: GestorArchivoBiblioteca ]
\end{verbatim}

\section{Ejecución del Programa}
A continuación se muestra la interfaz principal del programa durante su ejecución:

\begin{figure}[h]
    \centering
    % \includegraphics[width=0.6\textwidth]{interfaz_crud.png}
    \vspace{4cm} 
    \caption{Gestor de Biblioteca CRUD (reemplace la imagen)}
\end{figure}

\section{Solución: Lógica CRUD y Manejo de Archivos}
A continuación, se presenta un extracto fundamental para comprender cómo el código sobreescribe información en archivos secuenciales (método Update).

\begin{lstlisting}[caption={Clase GestorArchivoBiblioteca.java (Extracto)}]
// UPDATE
public static void actualizarLibro(String tituloBuscar, String nuevoTitulo, String nuevoAutor, String nuevoEstado) throws Exception {
    File file = new File(FILE_NAME);
    File temp = new File("temp.txt");
    boolean encontrado = false;

    // Abrimos el archivo original en lectura y un archivo temporal en escritura
    try (RandomAccessFile raf = new RandomAccessFile(file, "rw");
         RandomAccessFile tmpraf = new RandomAccessFile(temp, "rw")) {

        String line;
        while ((line = raf.readLine()) != null) {
            String[] datos = line.split("!");
            // Si es la linea que queremos actualizar, escribimos los datos nuevos
            if (datos.length == 3 && datos[0].equalsIgnoreCase(tituloBuscar)) {
                encontrado = true;
                String newRecord = nuevoTitulo + "!" + nuevoAutor + "!" + nuevoEstado;
                tmpraf.writeBytes(newRecord + System.lineSeparator());
            } else {
                // De lo contrario, copiamos la linea tal como estaba
                tmpraf.writeBytes(line + System.lineSeparator());
            }
        }
    }
    // Reemplazamos el archivo viejo por el nuevo modificado
    if (encontrado) {
        file.delete();
        temp.renameTo(file);
    }
}
\end{lstlisting}

\end{document}
